\section{Gaussian, exponential and atomic conditional mixtures}
\label{sec:mixturefamilies}
This section specializes the joint update to the distribution families of
primary interest. The same construction covers a Gaussian mixture, an
exponential mixture, either family with point masses, and a mixture containing
both continuous families. The normalized normal-means proposal is available
analytically; its hierarchical correction need not be conjugate.

\subsection{A normalized family with fixed supports}
For node \(u\) in cell \(i\), let \(h_{iu}=(P_{iu},X_i)\) collect its
permitted latent parents and fixed observed covariates. Suppress these
indices where unambiguous. Introduce a component label \(B\) with law
\begin{equation}\label{eq:familyprior}
\begin{split}
g(dv\mid h)={}&
\sum_{j\in\mathcal A}w_j(h)\delta_{c_j}(dv)
+\sum_{j\in\mathcal G}w_j(h)\Normal{v}{\mu_j(h)}{d_j^2(h)}\,dv\\
&+\sum_{j\in\mathcal E}w_j(h)\lambda_j(h)
             e^{-\lambda_j(h)v}\ind\{v>0\}\,dv .
\end{split}
\end{equation}
Here the three label sets are disjoint, \(w_j\geq0\), \(\sum_jw_j=1\),
\(d_j^2>0\), and \(\lambda_j>0\). Any set may be empty, provided at least
one component remains. A zero spike corresponds to a fixed \(c_j=0\).
The atoms and component supports do not depend on parents or trainable
parameters in this derivation. Moving an atom with a parent would introduce
deterministic support constraints and invalidate the simple density
cancellations below without further analysis.

The augmented conditional is \(g(dv,j\mid h)=w_j(h)G_j(dv\mid h)\).
Use Lebesgue measure on each continuous component and unit point mass on
each atomic component as the reference measure. Equivalently, a collapsed
prior uses Lebesgue measure plus the distinct fixed atoms. At an atom its
Radon--Nikodym derivative is the total atomic probability; continuous
densities evaluated at that point are not added to the atomic probability.
This convention removes any need to approximate a spike by a narrow Gaussian.

Neural networks may parameterize weights by softmax and positive scales or
rates by suitable positive transformations. A graph neural network is also
admissible if its latent inputs respect the specified directed acyclic graph.
If it introduces additional cross-cell dependencies, the child sets and
parallel update schedule must change accordingly. Fixed side information
may enter every conditional without being sampled. In all cases, each head
must specify a normalized probability distribution.

\subsection{Analytic normal-means proposals and moments}
\label{sec:familyproposal}
Condition on current parents, other latent variables and parameters, and
write the observation contribution as
\(\ell(v)=-av^2/2+bv\), with \(a\geq0\) and \(a=0\Rightarrow b=0\).
For every component define
\begin{equation}\label{eq:familyI}
I_j=\int e^{\ell(v)}G_j(dv),\qquad
Q_j(dv)=I_j^{-1}e^{\ell(v)}G_j(dv),\qquad
\widetilde w_j=\frac{w_jI_j}{\sum_r w_rI_r}.
\end{equation}
The exact own-prior normal-means posterior is
\(q_0(dv,j)=\widetilde w_jQ_j(dv)\).
The same unnormalized likelihood \(e^{\ell(v)}\) is used for every label;
one must not combine these \(I_j\) with normalized observation densities
for only some components.

\paragraph{Atomic components.}
For \(G_j=\delta_{c_j}\),
\begin{equation}\label{eq:familyatom}
\log I_j=-ac_j^2/2+bc_j,\qquad
Q_j=\delta_{c_j},\qquad
\E_{Q_j}V=c_j,\quad \E_{Q_j}V^2=c_j^2.
\end{equation}
In particular \(I_j=1\) for the zero spike. Its posterior probability still
changes through normalization against the continuous components.

\paragraph{Gaussian components.}
Let \(G_j=\N(\mu_j,d_j^2)\) and \(t_j=1+ad_j^2\). Completing the square gives
\begin{align}
V_j&=\frac{d_j^2}{t_j},&
M_j&=\frac{\mu_j+bd_j^2}{t_j},&
Q_j&=\N(M_j,V_j),\label{eq:familygauss}\\
\log I_j&=-\frac12\log t_j+
\frac{-a\mu_j^2/2+b\mu_j+b^2d_j^2/2}{t_j}.
\label{eq:familygaussI}
\end{align}
Thus \(\E_{Q_j}V=M_j\) and \(\E_{Q_j}V^2=M_j^2+V_j\).
The expression for \(\log I_j\) avoids subtracting large
\(\mu_j^2/d_j^2\) terms when the slab is narrow.

\paragraph{Exponential components.}
Let \(G_j(dv)=\lambda_j e^{-\lambda_jv}\ind\{v>0\}dv\).
For \(a>0\), put
\(z_j=(b-\lambda_j)/\sqrt a\) and \(m_j=(b-\lambda_j)/a\).
Since
\[
-av^2/2+(b-\lambda_j)v
=-a(v-m_j)^2/2+(b-\lambda_j)^2/(2a),
\]
integration over the positive half-line gives
\begin{align}
Q_j&=\N_{[0,\infty)}(m_j,a^{-1}),\label{eq:familyexp}\\
\log I_j&=\log\lambda_j+\tfrac12\log(2\pi/a)
                 +\tfrac12z_j^2+\log\Phi(z_j).\label{eq:familyexpI}
\end{align}
Writing \(r(z)=\varphi(z)/\Phi(z)\), its first two moments are
\begin{equation}\label{eq:familyexpmoments}
\E_{Q_j}V=\frac{z_j+r(z_j)}{\sqrt a},\qquad
\E_{Q_j}V^2=\frac{1+z_j^2+z_jr(z_j)}{a}.
\end{equation}
These follow either from truncated-normal moments or by differentiating
\(I_j\) with respect to \(b\). If \(a=b=0\), evaluate the limit directly:
\[
I_j=1,\qquad Q_j=\operatorname{Exp}(\lambda_j),\qquad
\E V=\lambda_j^{-1},\quad \E V^2=2\lambda_j^{-2}.
\]
This case occurs for a loading without a direct modality likelihood.
Its children can nevertheless inform its posterior.

For negative \(z\), an equivalent expression avoids cancellation in the
log normalizer:
\begin{equation}\label{eq:experfcx}
\log I_j=\log\lambda_j+\tfrac12\log\{\pi/(2a)\}
             +\log\operatorname{erfcx}(-z_j/\sqrt2),
\end{equation}
where \(\operatorname{erfcx}(x)=e^{x^2}\operatorname{erfc}(x)\).
Use log-CDF functions for the other branch and log-sum-exp for mixture
weights. The moment identities in \eqref{eq:familyexpmoments} can still
suffer cancellation for large negative \(z\); stable tail expansions or
specialized truncated-normal routines are needed in numerical solvers.
For example, if \(x=-z>0\), the unnormalized moments of
\(T=\sqrt a\,V\) are
\[
K_n(x)=\int_0^\infty t^n e^{-xt-t^2/2}\,dt
\sim x^{-n-1}\sum_{k=0}^\infty
\frac{(-1)^k(n+2k)!}{2^k k!\,x^{2k}}\qquad(x\to\infty).
\]
Use \(K_1/(\sqrt a K_0)\) and \(K_2/(aK_0)\) for the first two
moments, truncating this asymptotic expansion in the far tail.
The expansion follows by setting \(y=xt\), expanding
\(e^{-y^2/(2x^2)}\), and integrating each term against \(e^{-y}\).
It is an asymptotic expansion, not a convergent series at every \(x\).
There is no justification for replacing \(a=0\) by an arbitrary
positive precision.

A reflected exponential with fixed orientation \(\kappa_j\in\{-1,1\}\)
has density
\(\lambda_j e^{-\lambda_j\kappa_jv}\ind\{\kappa_jv>0\}\).
Apply \eqref{eq:familyexp}--\eqref{eq:familyexpmoments} to
\(T=\kappa_jV\) with \(b\) replaced by \(\kappa_jb\).
Multiply the first moment by \(\kappa_j\); leave the second moment unchanged.
The sign of a Gaussian mean, by contrast, does not restrict its support.

\paragraph{Mixture summaries.}
Without children, the first two posterior moments are
\[
\E_{q_0}V=\sum_j\widetilde w_j\E_{Q_j}V,\qquad
\E_{q_0}V^2=\sum_j\widetilde w_j\E_{Q_j}V^2.
\]
With a zero spike and otherwise continuous components,
\(\Pr(V=0)=\widetilde w_0\).
Under the convention
\(\mathrm{lfsr}=\min\{\Pr(V\leq0),\Pr(V\geq0)\}\),
the spike contributes to both probabilities; a positive exponential
contributes only to the second, and a Gaussian contributes its appropriate
normal-CDF probability. These are marginal posterior summaries, not
classification based on the largest component density.

\subsection{Exact hierarchical correction}
\label{sec:familycorrection}
During an augmented coordinate move keep each child's value \(V_c\) and
label \(B_c\) fixed. Define
\begin{equation}\label{eq:familyD}
D_u(v)=\sum_{c\in\ch(u)}
\log\{w_{cB_c}(h_c[v])G_{cB_c}(V_c\mid h_c[v])\}.
\end{equation}
The product is evaluated on the reference measures specified above.
For each component put
\[
J_{jr}=\int v^r e^{D_u(v)}Q_j(dv),\qquad r=0,1,2.
\]
Assuming a finite positive normalizer, the exact full conditional is
\begin{equation}\label{eq:familyfull}
p(dv,j\mid\mathrm{rest})
=\frac{w_jI_j}{\sum_h w_hI_hJ_{h0}}\,
       e^{D_u(v)}Q_j(dv).
\end{equation}
Its label probabilities are proportional to \(w_jI_jJ_{j0}\), and its
first two moments are obtained by replacing \(J_{j0}\) in the numerator
with \(J_{j1}\) and \(J_{j2}\), respectively. For an atom,
\(J_{jr}=c_j^r e^{D_u(c_j)}\), with \(c_j^0=1\).
This supplies a one-dimensional quadrature route with exact atomic masses.
The tilted continuous components generally cease to be Gaussian or
truncated Gaussian.

Alternatively draw \((v',j')\sim q_0\) and accept with probability
\begin{equation}\label{eq:familyMH}
\alpha=\min\{1,\exp[D_u(v')-D_u(v)]\}.
\end{equation}
To verify the rule, write the target as
\(\pi(v,j)=q_0(v,j)e^{D_u(v)}/Z_D\).
The forward accepted probability density is
\[
\pi(v,j)q_0(v',j')\alpha
=\frac{q_0(v,j)q_0(v',j')}{Z_D}
 \min\{e^{D_u(v)},e^{D_u(v')}\},
\]
which is symmetric in the two states. Rejections complete the transition
kernel, establishing detailed balance. Fixed-support mixture augmentation
therefore permits exact transitions between spikes and slabs without a
reversible-jump construction. This is the standard Metropolis principle
\cite{hastings}, specialized to analytic normal-means proposals.
It establishes invariance; irreducibility, adequate mixing and Monte Carlo
precision must be assessed separately.

For a Gaussian child component, include
\[
\log w_{cj}(h)-\tfrac12\log(2\pi d_{cj}^2(h))
-\frac{(V_c-\mu_{cj}(h))^2}{2d_{cj}^2(h)}.
\]
For a positive exponential child component, include
\begin{equation}\label{eq:expchild}
\log w_{cj}(h)+\log\lambda_{cj}(h)-\lambda_{cj}(h)V_c
\quad(V_c>0).
\end{equation}
For a fixed atomic child on its support, only the log weight remains.
The exponential rate normalizer and the Gaussian variance normalizer are
essential whenever they depend on the updated parent. In particular,
weights alone are insufficient for EMDN.

The obstruction to conjugacy is already present in the single-Gaussian
example \(V_c\mid v\sim\N(v^2,s^2)\), whose child log density contains
\(-(V_c-v^2)^2/(2s^2)\). Removing a spike does not remove this nonlinear
term. Conversely, linear means and constant variances in a Gaussian
loading DAG preserve Gaussian scalar conditionals, conditional on the
feature factors.

\subsection{Specialization to CGB, sharp CGB and two-slab CGB}
\label{sec:familycgb}
For the package's CGB prior,
\[
g(dv\mid h)=\{1-\pi(h)\}\delta_0(dv)
                +\pi(h)\N(dv;\mu,d^2).
\]
The normal-means masses are \(1-\pi\) and \(\pi I_G\), with
\(I_G\) from \eqref{eq:familygaussI}. The hierarchical slab log odds are
\begin{equation}\label{eq:cgbfamilyodds}
\log\frac{\Pr(B=1\mid\mathrm{rest})}{\Pr(B=0\mid\mathrm{rest})}
=\logit\pi+\log I_G+\log J_{G0}-D_u(0).
\end{equation}
Its conditional slab moments are \(J_{G1}/J_{G0}\) and
\(J_{G2}/J_{G0}\). If a CGB child has gate
\(\pi_c(v)=\sigmoid\eta_c(v)\), its variable log contribution is
\(B_c\eta_c(v)-\log\{1+\exp\eta_c(v)\}\).
The same expression applies to its sharp variant at fixed slab parameters.

Sharp CGB uses an effective slab variance \(d^2=\omega s^2\), with
\(\omega>0\), in the preceding formulas. A small positive variance is a
continuous slab, not a second point mass. If the base variance is freely
estimated, a fixed \(\omega\) merely reparameterizes the same family.
Persistent sharpening requires a fixed scale, an explicit constraint or
a specified regularizer. Repeatedly multiplying an unconstrained variance
estimate by \(\omega\) is not the unpenalized joint M-step.

The two-slab version has
\begin{equation}\label{eq:familycgb2}
g(dv\mid h)=\pi_0(h)\delta_0(dv)
+\pi_+(h)\N(dv;\mu_+,d_+^2)
+\pi_-(h)\N(dv;\mu_-,d_-^2),
\end{equation}
where \(\mu_+>0\), \(\mu_-<0\), and the weights sum to one.
The means are global parameters, and both Gaussians have support on
the whole real line. Its three posterior masses are proportional to
\[
\pi_0 e^{D_u(0)},\qquad
\pi_+I_+J_{+0},\qquad
\pi_-I_-J_{-0}.
\]
The moments follow by the same weighted integrals. With
\(d_\pm^2=\omega s_\pm^2\), this is the fixed-parameter model of
\texttt{cgb\_sharp\_2}. Setting \(\pi_0=0\) gives a two-Gaussian variant
without a spike; this is a distributional specialization, not an additional
registry name currently exposed by the package.

Setting the spike weight to zero in a one-slab CGB leaves a single global
Gaussian. Its gate disappears, so that node's own conditional prior no
longer depends on parents. Retaining parent dependence without a spike
requires multiple changing weights or parent-dependent means or variances,
as in EMDN. CASH and LCASH are further Gaussian-mixture specializations
with fixed zero-centered scale grids and their respective gate functions.

\subsection{Variational updates for the extended family}
\label{sec:familyvi}
For an alternative approximation, use the augmented mean-field family
\[
q=\prod_uq_u(V_u,B_u)\prod_{m,k}q^m_{F,k}
\]
from Section~\ref{sec:augmented}. Here \(a,b\) are the expected-likelihood
coefficients from Section~\ref{sec:vi}, rather than realized-value Gibbs
coefficients. Define \(D_u^q(v)\) as the expected log child contribution
under all other variational factors. Then
\begin{equation}\label{eq:familycavi}
q_u^*(dv,j)\ \propto\
\exp\{\ell(v)+\E_q\log[w_j(h)G_j(v\mid h)]+D_u^q(v)\}\,\nu_j(dv).
\end{equation}
All following expectations require finiteness. For a Gaussian component
this is the tilted Gaussian construction in
\eqref{eq:augweights}--\eqref{eq:augq}, with expected inverse variances
and expected precision-weighted means. Replacing these by the parameters
evaluated at mean parents is generally incorrect.

For an exponential component set
\begin{equation}\label{eq:expviAB}
B_j=b-\E_q\lambda_j(h),\qquad
C_j=\E_q\log w_j(h)+\E_q\log\lambda_j(h).
\end{equation}
Its unnormalized log coordinate density on \(v>0\) is exactly
\(-av^2/2+B_jv+C_j+D_u^q(v)\).
For \(a>0\), define
\begin{align}
H_j&=\sqrt{2\pi/a}\,
       \exp\{B_j^2/(2a)\}\Phi(B_j/\sqrt a),\\
Q_j^q&=\N_{[0,\infty)}(B_j/a,a^{-1}),&
J_{jr}^q&=\E_{Q_j^q}\{V^r e^{D_u^q(V)}\}.
\end{align}
The component mass is \(W_j=e^{C_j}H_jJ_{j0}^q\); conditional moments
are \(J_{j1}^q/J_{j0}^q\) and \(J_{j2}^q/J_{j0}^q\).
When \(a=b=0\), let \(\bar\lambda_j=\E_q\lambda_j(h)>0\).
Then \(H_j=1/\bar\lambda_j\) and
\(Q_j^q=\operatorname{Exp}(\bar\lambda_j)\), with the same tilted
integral formulas. Notice that \(e^{C_j}\) includes
\(\exp(\E\log\lambda_j)\); it cannot be replaced by
\(\E\lambda_j\) in the component mass.

For an atomic component at \(c_j\), the unnormalized mass is
\[
W_j=\exp\{\ell(c_j)+\E_q\log w_j(h)+D_u^q(c_j)\}.
\]
Normalize all Gaussian, exponential and atomic masses together.
For an exponential child component with
\(r_{cj}=q(B_c=j)\) and
\(m_{cj}=\E_q(V_c\mid B_c=j)\), its contribution to \(D_u^q(v)\) is
\begin{equation}\label{eq:expvichild}
r_{cj}\E_{P_c\setminus V_u}
\{\log w_{cj}(h_c[v])+\log\lambda_{cj}(h_c[v])
                    -\lambda_{cj}(h_c[v])m_{cj}\}.
\end{equation}
This factorization uses the stated mean-field independence between a
child node and its parents. A row-structured approximation instead retains
their joint expectation.

These are exact coordinate equations within the stated variational family,
provided the expectations and integrals are computed exactly. Monte Carlo
expectations, quadrature and parametric approximations to tilted components
are separate numerical approximations. The augmented and collapsed
mean-field problems are not interchangeable for overlapping mixtures.
An implemented variational algorithm must evaluate the corresponding
entropy and expected joint terms consistently.

\subsection{Learning the conditional distributions}
\label{sec:familymstep}
Let \(\mathcal Q_{\rm old}\) be either the joint latent posterior under
old parameters or a specified variational approximation. A node's
complete-data objective is
\begin{equation}\label{eq:familyQ}
\mathcal M_u(\theta_u)=
\sum_i\E_{\mathcal Q_{\rm old}}
\{\log w_{uB_{iu}}(h_{iu};\theta_u)
+\log G_{uB_{iu}}(V_{iu}\mid h_{iu};\theta_u)\}
+\mathcal R_u(\theta_u).
\end{equation}
Here \(\mathcal R_u\) is an explicitly specified regularizer or log
parameter prior; it is zero for the unpenalized empirical-Bayes objective.
The expectation is held fixed while optimizing \(\theta_u\).
Posterior samples must retain parents, values and labels from the same
joint draw. Latent posterior values are not new noisy observations for
another EBNM deconvolution step.

For global component parameters, write
\(n_j=\sum_i\E\ind\{B_{iu}=j\}\),
\(s_j=\sum_i\E[V_{iu}\ind\{B_{iu}=j\}]\), and
\(t_j=\sum_i\E[V_{iu}^2\ind\{B_{iu}=j\}]\).
Without regularization and with positive occupancy, Gaussian components
have the unconstrained updates
\begin{equation}\label{eq:familygaussmstep}
\mu_j^{\rm new}=s_j/n_j,\qquad
(d_j^2)^{\rm new}=t_j/n_j-(s_j/n_j)^2.
\end{equation}
If the mean is constrained to a closed interval \(\mathcal I_j\), project
\(s_j/n_j\) onto that interval and use
\(t_j/n_j-2\mu_j^{\rm new}s_j/n_j+(\mu_j^{\rm new})^2\)
for the variance. A strict sign constraint can have only a boundary
supremum; using \([\epsilon,\infty)\) or \((-\infty,-\epsilon]\)
with fixed \(\epsilon>0\) specifies an attained constrained problem.
Variance bounds are enforced by projection onto the specified positive
interval. For a global exponential rate,
\begin{equation}\label{eq:familyexpmstep}
\lambda_j^{\rm new}=n_j/s_j
\end{equation}
when \(s_j>0\); bounds or regularization are required at degeneracies.
Unoccupied components do not have identified parameter updates.
For global weights, \(w_j^{\rm new}=n_j/N\).

Neural conditional parameters generally require numerical optimization of
\eqref{eq:familyQ}. For a Gaussian draw their gradients with respect to
\(\mu\) and \(\log d\) are
\[
\frac{v-\mu}{d^2},\qquad \frac{(v-\mu)^2}{d^2}-1;
\]
for an exponential draw the gradient with respect to \(\log\lambda\)
is \(1-\lambda v\). A softmax logit receives
\(\ind\{B=j\}-w_j(h)\). The network chain rule then accounts for its
chosen parameterization, including shared hidden layers. These gradients
describe complete-data learning, not the derivative of an isolated
normal-means marginal likelihood.

Exact EM increases the observed-data likelihood when its E-step is exact
and its M-step increases the corresponding objective. Finite Monte Carlo
EM need not do so, even if it improves the fixed-sample objective.
Variational coordinate ascent instead concerns its particular ELBO.
Mixture variance collapse, rate divergence and overfitting by flexible
heads motivate explicit admissible parameter sets and held-out tuning;
they are not resolved by the algebra of a scalar proposal.

\subsection{Scope of the derivation and implementation}
The formulas above supply exact normal-means proposals and invariant
hierarchical coordinate kernels for the stated families. They do not claim
an analytic general full conditional or uniformly fast mixing.
An exponential mixture on the positive half-line has a decreasing
continuous density:
\[
\frac{d}{dv}\sum_jw_j\lambda_je^{-\lambda_jv}
=-\sum_jw_j\lambda_j^2e^{-\lambda_jv}<0\quad(v>0).
\]
Consequently, it is not a universal family for positive conditional
densities and cannot represent a continuous mode away from zero.
Adding a zero spike does not change this restriction. The choice between
Gaussian and exponential components therefore changes both support and
shape assumptions, rather than only computational convenience.
Their HMM counterpart remains a sequence block update
(Section~\ref{sec:hmm}); the loading hierarchy does not introduce children
of \(F\) when feature factors are not inputs to the loading priors.

The experimental code currently implements Gaussian and fixed-atom
conditional mixtures across the eight scalar learned registry families.
Exponential-mixture inference and its variational counterpart in this
section are derived extensions, not newly implemented solver options.
The independent verification script described in Section~\ref{sec:checks}
tests their local identities and the TwinEB factorization below.
